% ============================================================
% Chapter 11: The Chinese AI Startup Ecosystem — A Deep Analysis
% Book: The AI Startup Landscape (2022--2026)
% Language: Simplified Chinese (per book convention for China chapter)
% ============================================================

\chapter{中国 AI 创业生态深度解析}
\label{ch:china}

% China's parallel AI ecosystem — investors, big tech labs,
% independents, regulation, and the US-China tech competition.
% This chapter provides the most comprehensive survey of China's
% AI startup landscape as of March 2026.

如果说硅谷是全球 AI 创业的"第一极"，那么中国就是无可争议的"第二极"——
而且这个"第二极"正在以自己独特的方式重新定义 AI 创业的规则。

2025 年，中国人工智能一级市场投资事件共计 1579 起，同比激增 75\%，创下近五年新高；
投资总额达 1504.27 亿元人民币，同比增长 24\%（IT 桔子，2026 年 1 月）。
进入 2026 年一季度，热度不减反增：仅前三个月即披露 88 起融资事件、
合计金额超 200 亿元，银河通用单轮 25 亿元刷新具身智能融资纪录。
与此同时，智谱和 MiniMax 于 2026 年 1 月 8--9 日先后在港交所挂牌上市，
标志着中国大模型公司正式进入公开资本市场的"成年礼"。

但中国 AI 创业生态的独特性，远不止资金量级。它是一个由国家政策引导、
大厂生态笼罩、独立创业公司突围、开源社区反哺的四层结构，
每一层都有自己的逻辑，每一层之间又存在复杂的竞合关系。
理解这个生态，需要同时掌握资本、技术、政策和文化四种语言。

\begin{corebox}[中国 AI 创业生态的四层结构]
\textbf{第一层：大厂生态层。}字节跳动（豆包/Seed）、腾讯（混元/元宝）、阿里巴巴（千问/通义）、百度（文心）构成中国 AI 的"基础设施"。它们既是模型提供者，也是最大的 AI 应用分发渠道，更是创业公司绕不过去的生态入口。2026 年春节期间，豆包、元宝、千问三款应用的营销战总花费超百亿元，共同跻身"月活亿级俱乐部"。

\textbf{第二层：独角兽层。}智谱（GLM 系列）、MiniMax（海螺 AI/Talkie）、月之暗面（Kimi）、DeepSeek、阶跃星辰等"AI 六小龙"在基座模型赛道形成寡头格局。其中智谱和 MiniMax 已上市，月之暗面 2026 年 3 月估值飙升至 180 亿美元（约 1200 亿元），阶跃星辰完成 50 亿元 B+ 轮融资。

\textbf{第三层：垂直创业层。}覆盖 AI 办公、AI 教育、AI 电商、AI 社交、AI 搜索、具身智能等数十个垂直赛道。2025 年标签含"AI 应用"且获得融资的公司达 930 家，融资总额 1070.7 亿元——平均每天 2.6 家公司拿到融资。

\textbf{第四层：开源社区层。}以 DeepSeek、千问（Qwen）为代表的中国开源模型在 Hugging Face 上的下载量和衍生模型数已跻身全球前列。千问开源模型超 400 个，衍生模型突破 20 万个，下载量超 10 亿次。DeepSeek 则以"低成本高性能"路线引发全球震动。
\end{corebox}

本章将逐层拆解这个复杂生态。我们从资本入手（§\ref{sec:china-investors}），
依次分析大厂 AI 内部创业（§\ref{sec:china-bigtech}）、独立创业公司矩阵
（§\ref{sec:china-independents}）、政策与监管环境（§\ref{sec:china-policy}）、
中美竞争态势（§\ref{sec:china-competition}）、媒体生态（§\ref{sec:china-media}），
最后以四城产业集聚的比较（§\ref{sec:china-clusters}）收束全章。


% ============================================================
%  §11.1  PRIMARY MARKET LANDSCAPE
% ============================================================

\section{一级市场全景}
\label{sec:china-investors}

2025 年的中国 AI 一级市场，呈现出一个清晰的结构性特征：
\textbf{资金向头部集中，阶段向早期前移}。
IT 桔子数据显示，全年 1579 起投资事件中，早期投资（种子轮至 Pre-A 轮）高达 754 起，
同比涨幅 90\%，占比过半。天使轮和 A 轮各占近三分之一，C 轮及以后仅占八分之一。
这意味着 VC 们的策略已从"造王者"式的押注后期巨头，转向"广撒网"式的早期布局。

与此同时，头部项目的单轮融资金额在持续刷新纪录。
2025 年共有 6 家公司融资超 10 亿元，其中银河通用机器人累计融资超 32 亿元。
2026 年一季度，阶跃星辰 50 亿元 B+ 轮成为国内 AI 领域最大单轮融资之一，
月之暗面则在不到 90 天内完成三轮融资、合计超 22 亿美元，估值从 43 亿美元飙至 180 亿美元。

\subsection{头部 VC 的 AI 布局}
\label{subsec:china-vc-landscape}

\subsubsection{红杉中国（Hongshan）：从 xbench 到学术高地}

红杉中国（2024 年更名为 Hongshan）是中国 AI 创业投资中最具系统性的玩家。
其布局横跨基座模型（月之暗面、智谱、MiniMax）、具身智能（自变量机器人、千寻智能）、
AI 芯片（寒武纪）等全产业链。

红杉中国的独特之处在于其"学术高地"策略。2025 年 5 月，红杉中国推出 xbench——
国内首个由投资机构发起的 AI 基准测试平台，采用"双轨评估体系"（理论上限 + 真实场景效用）
和"长青评估机制"（持续动态更新），在 AI 行业引发巨大反响。
随后开源了 xbench-ScienceQA 与 xbench-DeepSearch 两个评测集。
2026 年 1 月，xbench 更在两周内连发 BabyVision（多模态理解）和 AgentIF-OneDay（长程 Agent 评估）两个新评测集，
同期发布四篇学术论文。这种"用学术影响力吸引 Founder"的打法，
让红杉中国在 AI 投资的"抢人大战"中建立了独特的护城河。

2026 年 2 月，红杉中国领投千寻智能近 20 亿元融资，使其估值突破百亿，
成为具身智能领域又一百亿独角兽。在 AI 投资人口中流传的一句话——"想接触的 Founder 对我们爱搭不理"——
恰恰说明了 AI 时代 VC 与创始人权力关系的逆转。红杉中国的应对之策，
是通过 xbench、学术论文和评测生态，将自己打造成 AI 开发者社区的"学术合伙人"。

\subsubsection{高瓴资本：500 亿募资的表层逻辑}

高瓴资本在 AI 领域的投资同样布局广泛。2025 年 12 月，外媒消息称高瓴启动新一轮
约 70 亿美元（约 500 亿元）的美元 PE 基金募资，若规模确实将成为 2022 年以来
亚洲本土机构最大的单只 PE 基金。

高瓴在 AI 领域的投资逻辑经历了一次重要转变。早期阶段，高瓴以"广撒网"方式
参投了多家大模型公司（包括百川智能、MiniMax 等），但随着"百模大战"格局明朗，
高瓴逐步将重心转向\textbf{应用层和产业 AI}。
红杉与高瓴联合押注讯兔科技等项目，体现了头部 VC 在 AI 赛道"抱团取暖"的趋势。

高瓴的优势在于其全周期投资能力——从一级市场的 VC/PE 到二级市场的对冲基金，
这使其能够在大模型公司从融资到上市的全过程中持续参与。
MiniMax 上市后，高瓴持股 1.11\%，虽非大股东，但其"跨周期陪跑"的能力
使其在 AI 创业者中保持了独特的号召力。

\subsubsection{启明创投与五源资本：早期猎手}

启明创投是阶跃星辰 50 亿元 B+ 轮的投资方之一，其在 AI 领域的布局以"投早投小"为核心策略，
聚焦于具备技术差异化的早期项目。五源资本同为阶跃星辰的投资方，
其创始合伙人刘芹长期关注 AI 基础设施和底层技术方向。

\subsubsection{奇绩创坛与真格基金：技术创始人的"第一通电话"}

陆奇创立的奇绩创坛（MiraclePlus）延续了 Y Combinator 的孵化模式，
是中国 AI 早期创业者最重要的"起跑线"之一。其"Demo Day"已成为中国 AI 创投界的标志性事件。

真格基金则是月之暗面天使轮的重要投资方。徐小平团队对"技术型创始人"的偏好，
使其在 AI 浪潮早期就布局了一批核心标的。月之暗面从天使轮到 180 亿美元估值的
跃升，为真格基金带来了丰厚的账面回报。

\subsection{资本结构的三大变化}
\label{subsec:china-capital-shifts}

\subsubsection{国资入场：从陪跑到主导}

2025 年中国 AI 投资最显著的结构性变化，是\textbf{国有资本从陪跑席走向牌桌中心}。
智谱 IPO 的基石投资者包括北京核心国资、头部保险资金及大型公募基金，
合计认购 29.8 亿港元。杭州城投产业基金和上城资本以超 10 亿元战略投资智谱，
反映出地方政府通过资本纽带加速布局大模型产业基础设施的意图。

阶跃星辰 50 亿元 B+ 轮的投资方中，上海国投、国寿股权、浦东创投、徐汇资本赫然在列，
国资占据了核心位置。更引人注目的是，中国移动以 3 亿美元领投银河通用机器人，
央企的直接入场标志着 AI 投资已上升至国家战略层面。

北京机器人产业发展投资基金、深创投等国资机构在具身智能领域频频出手，
深创投 2025 年布局了 40 个机器人全产业链项目。国资入场不仅带来了资金，
更带来了政策协调能力和产业落地资源——这是纯财务投资者无法提供的。

\subsubsection{产业资本跨界：从投资者到"客户+股东"双角色}

阿里巴巴是 2025 年具身智能赛道投资总额排名第一的投资方，连投四家具身智能企业。
宁德时代以"客户+股东"双角色绑定银河通用——千台订单反哺被投企业，
形成了"我投你、你买我、一起赚钱"的产业闭环。

美团、京东投资机器人瞄准仓储物流场景，北汽、吉利瞄准汽车产线场景。
产业资本的深度参与，使得中国 AI 投资从"讲故事、画大饼"的模式
转向"投完就能用、用完能验证"的产业协同模式。

\subsubsection{中东资本进场}

阿布扎比投资局（ADIA）出现在 MiniMax 港股基石投资者名单中，
认购 6500 万美元。中东主权基金还投资了三家中国机器人公司。
在中美科技竞争的大背景下，中东资本的进入为中国 AI 企业提供了
"第三条融资渠道"——既不依赖美元 VC，也不完全依赖国资。

\subsection{退出通道：IPO 潮定义周期分水岭}
\label{subsec:china-exit}

2026 年 1 月 8--9 日，智谱与 MiniMax 前后脚登陆港交所，掀开了中国大模型公司
IPO 潮的序幕。

\textbf{智谱}（SEHK: 02513）以 116.2 港元发行价上市，募资 43.48 亿港元，
上市首日收盘市值 579 亿港元，后迅速攀升至约 1000 亿港元。
公开发售获 1159.46 倍超额认购，国际发售获 15.28 倍认购。
智谱是港股首家以通用大模型（AGI 基座模型）为核心资产的上市企业，
其"双轮驱动"模式以高毛利的本地化部署为基石（2025 上半年贡献 85\% 收入，毛利率 59\%），
云端 API 服务为增长引擎。

\textbf{MiniMax}（SEHK: 00100）以 165 港元的定价上限发行，募资 48.2 亿港元，
上市首日暴涨一度超 90\%，收盘市值突破 900 亿港元。
公开发售超额认购 1837 倍，国际发售认购需求达 300 亿美元，
创下近年来港股 IPO 机构认购历史纪录。MiniMax 的差异化在于其"B+C 双轮驱动"模式
和海外收入占比超 70\%——Talkie/海螺 AI 触达全球超 200 个国家的 2.12 亿用户，
使其成为国产大模型全球化的标杆。

同期上市的还有 AI 芯片公司天数智芯（SEHK: 9903）和壁仞科技（SEHK: 6082），
加上此前已上市的 GPU 四小龙中的摩尔线程和沐曦（详见第二章），
中国 AI 产业的"集体上市潮"覆盖了从芯片到模型的完整产业链。

\begin{warnbox}[IPO 热潮背后的估值泡沫风险]
MiniMax 上市后市值一度突破 3826 亿港元，超越百度（3322 亿港元）、携程、快手等
老牌互联网巨头。然而，截至 2025 年前 9 个月，MiniMax 亏损超 5 亿美元，
约为收入的近 10 倍。智谱同样处于大幅亏损状态。

市场对大模型公司的估值逻辑建立在"AI 是下一代通用技术平台"的信念之上——
投资者买入的不是当期利润，而是未来十年 AI 渗透所有行业的增量空间。
但当所有人都在为同一个梦想付费时，泡沫风险不容忽视。

一个值得思考的问题是：智谱和 MiniMax 的合计市值已超过万亿港元，
而它们 2025 年的合计营收可能不到 50 亿元——
这意味着隐含市销率超过 200 倍。即便按最乐观的增速假设，
投资者也需要数年时间才能看到估值与基本面的匹配。
\end{warnbox}


% ============================================================
%  §11.2  BIG TECH AI INTRAPRENEURSHIP
% ============================================================

\section{大厂 AI 内部创业}
\label{sec:china-bigtech}

中国 AI 创业生态的独特之处在于：大厂不仅是投资者和基础设施提供者，
更是 AI 产品层面最凶猛的竞争者。字节跳动、腾讯、阿里巴巴、百度——
四家公司各自投入数百亿元，用截然不同的策略抢占 AI 时代的入口。

\subsection{字节跳动：豆包与 Seed 团队的"All In"时刻}
\label{subsec:china-bytedance}

豆包是中国 AI 应用史上最重要的产品之一。2025 年 12 月 24 日，
36 氪独家报道豆包日活跃用户（DAU）突破 1 亿大关——
这是中国移动互联网历史上首个 DAU 破亿的 AI 原生应用，
也使其成为字节跳动继今日头条、抖音之后的第三款亿级 DAU 产品。

更令人惊叹的是达到这一里程碑的效率。据字节内部人士透露，
豆包的用户增长（UG）和市场推广费用，是字节历史上所有破亿 DAU 产品中花费最低的。
豆包 Seedream 生图模型和 Seedance 生视频模型持续激活创作场景，
三宫格等 P 图玩法连续登上热搜，每天带动的自然下载量高达百万级。

截至 2026 年 2 月，QuestMobile 数据显示豆包月均活跃用户达 2.3 亿，
连续两个季度登顶行业榜首，月均下载用户数连续三季度位居行业第一。
火山引擎宣布豆包大模型日均 Token 调用量突破 50 万亿，较上年增长超 10 倍。

字节跳动的 AI 战略展现出极强的\textbf{垂直整合能力}：

\begin{itemize}
  \item \textbf{基础设施层（MaaS）：}火山引擎以"极致性价比"抢占 B 端市场，
    2025 年上半年在 MaaS 层外部客户 Tokens 调用量市场份额位居第一（49.2\%）；
  \item \textbf{模型层：}豆包大模型（Doubao-Pro/Lite）、Seed 1.8 系列，
    以稀疏混合专家（MoE）架构为核心，定价比 GPT-4 便宜 50 倍、比行业均价低 62.7\%；
    2025 年 6 月发布的豆包 1.6 系列模型（flash/thinking/all-in-one 三版本）支持 256K 超长上下文，
    首创按"输入长度"区间定价策略，综合使用成本降至 1.5 版本的三分之一；
  \item \textbf{应用层：}豆包 APP 为核心 C 端入口，Coze（扣子）为低代码 Agent 开发平台，
    豆包手机（努比亚 M153）完成"硬件+AI"布局，接入抖音电商打通"AI+消费"场景。
\end{itemize}

\subsubsection{产品矩阵：从单一助手到"App 工厂"}

字节围绕豆包大模型构建了业内最庞大的 AI 产品矩阵，
延续了其在移动互联网时代"App 工厂"的产品基因：

\begin{itemize}
  \item \textbf{豆包（对话助手）：}C 端旗舰，DAU 破亿，覆盖问答、创作、翻译、代码等通用场景；
  \item \textbf{豆包 AI 搜索：}在豆包 APP 内集成深度搜索能力，
    直接对标秘塔和 Perplexity，通过抖音信息流和搜索数据构建差异化信息源；
  \item \textbf{扣子/Coze：}面向开发者和企业的 AI Agent 开发平台，
    支持无代码/低代码方式构建智能体，已接入国内外多个大模型，
    定位为"AI 时代的飞书"——从办公协作升级为 Agent 编排平台；
  \item \textbf{即梦 AI：}专注 AI 创作与设计，整合 Seedream 文生图模型和 Seedance 视频生成模型，
    以"三宫格"等创意 P 图玩法频登热搜，日均带动百万级自然下载；
  \item \textbf{剪映 AI：}中国用户量最大的 AI 视频编辑工具，
    集成"AI 一键成片""智能字幕""AI 换脸"等功能，
    已成为抖音内容生态的核心创作基础设施；
  \item \textbf{豆包爱学：}聚焦 K12 学科教育，凭借抖音流量低成本导流快速起量；
  \item \textbf{猫箱（Cat Box/星野）：}AI 角色扮演与社交陪伴应用，
    基于豆包角色扮演模型，主打虚拟角色互动和情感陪伴。
\end{itemize}

\subsubsection{海外版图：Cici/Dola 与 Dreamina}

字节的 AI 出海战略同样激进。豆包海外版最初以"Cici"品牌发布，
2025 年下半年在印尼、马来西亚、菲律宾、墨西哥和英国等市场始终位列
Google Play 免费应用下载榜前 20 名。在墨西哥，Cici 曾连续一周蝉联
Google Play 下载量第一。Sensor Tower 数据显示其下载量持续保持高位。

2025 年 12 月底，Cici 更名为 Dola 并加速投放，日活跃用户突破千万。
同期，即梦海外版 Dreamina AI 启动"新年攻势"加大投放力度；
字节还面向海外市场上线了对标 Manus 的工作场景 Agent 产品 AnyGen。
C 端与 B 端的双线布局，使字节具备了在海外 AI 市场参与体系化竞争的能力。

值得注意的是，Cici/Dola 的产品页面刻意淡化了与字节跳动的关联——
这一策略与 TikTok 早期的品牌独立路线如出一辙，
目的是在地缘政治敏感的海外市场降低"中国 App"的标签效应。

\subsubsection{组织与投入}

在组织层面，2025 年字节完成了 AI Lab 等多团队的整合，由前 Google Brain
研究员吴永辉出任 Seed 团队"一号位"，全面聚焦大模型基础研究与应用落地。
为激励团队，字节为 Seed 成员提供额外期权奖励，并取消部分长期研究项目的季度 OKR 考核——
这在以"效率至上"著称的字节内部极为罕见，折射出公司对 AI 研发的战略重视程度。
创始人张一鸣亲自下场参与 AI 战略制定，资本开支持续加码，
高薪从 Google、Meta 等公司招募 AI 顶尖人才，体现出"All In AI"的决心。

火山引擎成为 2026 年央视春晚独家 AI 云合作伙伴，豆包配合上线多种互动玩法——
十亿观众规模的春晚，正如 2015 年微信红包之于移动支付，
字节试图用这一"国民级舞台"将豆包推向真正的国民级应用。
豆包大模型已接入字节内部 50 余个业务线，包括抖音、今日头条、番茄小说、
飞书、巨量引擎等核心产品，日均 Token 调用量突破 50 万亿。

\subsection{腾讯：混元大模型与元宝的"后发先至"}
\label{subsec:china-tencent}

腾讯在大模型赛道的起步略晚于字节，但凭借微信生态的天然优势实现了快速追赶。

2026 年 2 月 18 日，腾讯宣布元宝 DAU 突破 5000 万，MAU 达到 1.14 亿，
跻身国内 AI 原生应用 Top 3。春节期间"上元宝，分 10 亿"活动中，
元宝主会场累计抽奖超 36 亿次，用户通过"创作"栏完成 AI 任务超 10 亿次。
随后 MAU 进一步增长至超过 1 亿（2026 年 3 月财报数据）。

腾讯的 AI 战略围绕三个关键词展开：\textbf{模型升级、产品矩阵、生态嵌入}。

\textbf{模型升级：}混元大模型持续迭代——2025 年 12 月上线混元 2.0，
2026 年 1 月混元图像 3.0 接入元宝实现一句话 P 图（春节期间 AI 生图日均调用量增长 30 倍），
混元 3.0 计划于 2026 年 4 月向外开放。
混元 3D 系列模型在开源社区的下载量突破 300 万。

\textbf{产品矩阵：}腾讯多款 AI 产品密集升级——
元宝一年内累计更新超百次；AI 工作台 ima 月活超 1300 万，知识库文件超 4.2 亿；
QQ 浏览器 AI 累计服务超 1.3 亿用户；搜狗输入法 AI 用户破 1 亿，移动端月活 6.7 亿。

\textbf{生态嵌入：}腾讯最大的护城河是微信。公众号、视频号内容构成元宝底层数据的
重要来源，元宝 AI 能力已接入 QQ 音乐、腾讯会议、微信公众号评论区等数十款核心应用。
"元宝派"AI 社交新形态让用户与好友共同邀请 AI 参与聊天，
以及"小龙虾"系列 Agent（WorkBuddy、QClaw）在微信聊天界面直接调度——
这些举措使腾讯的 AI 能力深度嵌入了用户最高频的社交和工作场景。

在组织层面，腾讯 2025 年 12 月新设 AI Infra 部、AI Data 部和数据计算平台部，
进一步强化 AI 研发体系。同时加大原生 AI 人才引进，
以年轻化力量重构组织，升级大模型研发架构。
2025 年全年资本开支 792 亿元、研发投入 857.5 亿元，
均创历史新高。腾讯 2025 年全年营收 7517.7 亿元（+14\%），
经营利润 2806.6 亿元（+18\%），AI 已在广告、游戏和社交等核心业务中产生实质性增益。

\subsubsection{微信智能体：最高优先级的绝密项目}

2026 年 3 月，《财经》杂志报道腾讯正在为微信开发一款 AI 智能体，
该智能体将与微信内数百万个小程序连接，提供从打车到网购等各类服务。
腾讯至少从 2025 年上半年就开始筹办该项目，并将其视为"最高优先级的绝密项目"。

微信（月活 14.14 亿，中国排名第一的超级 App）若成功接入 AI Agent，
将创造全球最大的"AI + 超级 App"生态——这相当于 ChatGPT 内嵌于 WhatsApp，
同时直接调度支付、电商、出行、生活服务的全场景能力。
然而，接近腾讯人士表示，因涉及隐私安全，该项目的发布时间仍是未知数。
微信将用户体验放在首位，增加任何新功能都要确保不影响用户体验；
同时，小程序背后的厂商是否愿意向微信智能体开放权限，也是关键未知数。

\subsubsection{"小龙虾"Agent 矩阵}

腾讯在 Agent 层面的布局以"小龙虾"（OpenClaw 浪潮的腾讯版本）为代表，
推出了 WorkBuddy（工作场景 Agent）和 QClaw（通用 Agent）等产品，
可在微信聊天界面直接调度各类任务。
这种将 AI Agent 深度嵌入用户最高频社交场景的策略，
使腾讯在 Agent 时代具备了独特的入口优势——
用户不需要打开新 App，在微信对话框中即可让 AI 代劳。

马化腾在 2025 年全年财报中表示："我们的混元 3.0 大语言模型智能水平持续提升，
元宝、WorkBuddy 及 QClaw 等 AI 产品产生实际效用，
这些令人鼓舞的初期迹象表明 AI 投入将为我们开拓新的机遇。"

\subsection{阿里巴巴：千问品牌统一与开源野心}
\label{subsec:china-alibaba}

2026 年 2--3 月，阿里巴巴完成了一次重要的品牌整合——
将旗下 AI 产品体系名称统一为"千问"：
"千问大模型"为整体品牌中文名（英文 Qwen），"千问 APP"为 C 端旗舰应用，
"通义实验室"为组织名称。这一整合结束了此前"通义千问""千问""Qwen"
多名称混用导致的品牌混乱。

千问 APP 自 2025 年 11 月公测以来不到两个月 MAU 即突破 1 亿，
"一句话下单"功能直接调用淘宝商品信息，春节期间用户完成近 2 亿次下单。
QuestMobile 数据显示，春节期间千问 DAU 达到 7352 万，增幅 940\%。

但阿里 AI 战略的真正杀手锏是\textbf{开源}。截至 2026 年 3 月，
千问开源模型超过 400 个，衍生模型突破 20 万个，下载量超 10 亿次——
不仅是中国开源模型数量第一，在全球小模型的覆盖度和生态成熟度上同样位居前列。
2026 年除夕开源发布的千问 3.5 在 Hugging Face 上包揽趋势榜前 8 位，
马斯克公开点赞"智能密度令人印象深刻"。

然而，阿里千问也面临挑战。2026 年 3 月 4 日，千问团队技术负责人林俊旸离职，
引发行业震动——这位阿里最年轻的 P10 级技术负责人是将千问从零打造为全球顶级
开源生态的核心灵魂。阿里 CEO 吴泳铭随即宣布成立"基础模型支持小组"，
强调将继续坚持开源策略并加大 AI 研发投入。
千问旗舰模型 Qwen3.5 在 LMArena 榜单上位列第 18 位，
未能复现此前 Qwen3-Max Preview 进入前三的辉煌。

阿里在 AI 基建上的投入同样巨大：2025 年 2 月宣布未来三年投入 3800 亿元用于
AI 基础设施建设，旗下平头哥（PPU 芯片）传出独立 IPO 消息。
2026 年 3 月财报电话会上，阿里披露平头哥 GPU 芯片已累计规模化交付 47 万片，
并宣布 AI 战略商业目标：未来五年云和 AI 商业化收入突破 1000 亿美元。

\subsubsection{阿里云百炼：模型即服务的超级平台}

阿里云百炼（Bailian）平台已成为中国最大的大模型 API 服务平台之一，
上线 100 多款国内外主流模型，提供 400 多个 AI Agent 模板与服务。
2026 年 2 月，百炼推出"最强 Coding Plan"——
包含千问 3.5、GLM-5、MiniMax M2.5、Kimi K2.5 四大顶尖开源模型的 API 订阅套餐，
用户可实现多模型自由切换。全球云厂商中，仅阿里云提供此类多模型订阅服务。
Pro 套餐新用户首月仅需 39.9 元，可享每月 90000 次请求，
大幅降低了高频次 AI 编码场景的成本。

此外，阿里还推出 Qwen Code——一款 AI 编程助手（CLI 与 IDE 插件），
可在 Claude Code、Cline、OpenClaw 等 AI 工具上无缝使用百炼模型，
标志着阿里在开发者工具链上与 OpenAI、Anthropic 直接竞争。

\subsubsection{钉钉 AI：从协作平台到"Agent OS"}

钉钉在无招回归后大刀阔斧重塑产品。2025 年推出录音卡片 DingTalk A1
和 AI 终端 DingTalk Real 两款硬件，以"Agent OS"构建 AI 时代的工作操作系统。
钉钉接入通义千问能力，服务数百万企业客户，
其策略从叶军时代的纯平台化（PaaS only）转回"软硬一体"，
硬件 + AI Agent 成为其差异化飞书的关键。
2026 年奥运会期间，阿里宣布千问成为米兰冬奥官方大模型，
进一步提升了千问的全球品牌认知度。

\subsection{百度：文心大模型与搜索转型}
\label{subsec:china-baidu}

百度是中国最早全面拥抱大模型的互联网巨头。文心一言（ERNIE Bot）于 2023 年 3 月
上线，是国内首批向公众开放的大模型应用。2025 年 11 月，文心 5.0 正式版上线——
这是原生全模态大模型，参数达 2.4 万亿，采用原生全模态统一建模技术，
支持文本、图像、音频、视频等多种信息的输入与输出。

\subsubsection{财务全景：AI 业务占比突破 43\%}

2026 年 2 月发布的 2025 年全年财报揭示了百度 AI 转型的实质性进展。
2025 年总营收 1291 亿元（同比 $-3\%$），传统广告业务的下滑仍在拖累整体表现，
但\textbf{核心 AI 新业务全年收入达 400 亿元}（同比 $+48.1\%$），
四季度 AI 业务收入占百度一般性业务收入的 43\%，超出市场预期。

具体拆分：智能云基础设施营收 198 亿元（$+33.8\%$），
AI 应用营收 102 亿元，AI 原生营销服务营收 98 亿元（$+308.3\%$）。
AI 高性能计算设施的订阅收入四季度同比增长 143\%，较三季度 128\% 进一步加速。
百度智能云在 2025 年全年大模型相关招标项目中，
中标项目数及中标金额均居行业首位，连续两年蟬联"标王"。

李彦宏在财报电话会上强调"应用比模型更重要"，
明确百度将坚持应用驱动的发展路径，以真实需求反哺和迭代模型能力。
这一论断既是对 DeepSeek "低成本高性能"路线的回应，
也反映出百度在基座模型竞赛中策略性后撤、转向应用侧的务实判断。

\subsubsection{萝卜快跑：全球 Robotaxi 龙头}

萝卜快跑（Apollo Go）是百度 AI 战略中最具独立增长潜力的业务。
2025 年四季度，全球无人驾驶出行服务次数达 340 万次，同比增长超 200\%，
季度内每周出行次数峰值超 30 万。截至 2026 年 2 月，
萝卜快跑累计提供全球出行服务次数超 2000 万次，足迹覆盖全球 26 个城市。
自动驾驶总里程累计超 3 亿公里，其中全无人驾驶里程超 1.9 亿公里。

在商业化方面，受益于第六代车款 RT6 成本低于 3 万美元的优势，
萝卜快跑已在武汉达成单位经济效益损益平衡。
国际化步伐同样加速——2025 年 11 月获得阿布扎比首批全无人驾驶商业化运营许可，
计划 2026 年进军伦敦、瑞士圣加仑、迪拜、首尔及香港等城市。

\subsubsection{文心助手与搜索转型}

百度的差异化在于其\textbf{搜索基因}。作为中国最大的搜索引擎，
百度将 AI 深度融入搜索场景——"百度搜索接入深度思考模型"成为 2025 年 AI 搜索赛道
的标志性事件。2025 年 12 月，百度 App 月活用户 6.79 亿，
文心助手（原文心一言）月活用户达 2.02 亿。
春节红包活动启动以来，文心助手月活跃用户同比增长 4 倍。

百度还成立了个人超级智能事业群（PSIG），
将百度文库与百度网盘整合，加速推动 AI 应用创新。
慧播星数字人开播规模 2025 年同比增长 202\%，收入同比增长 228\%，
已赋能 30 多个行业。无代码生成平台"秒哒"在 IDC 评测中表现出行业领先性能。

昆仑芯（从百度分拆，详见第二章）于 2026 年 1 月向港交所递交 IPO 申请，
目标融资 20 亿美元，进一步扩充百度系 AI 芯片的资本实力。

但百度面临的挑战也最为严峻：在 C 端，文心助手的用户量被豆包、元宝、千问全面超越；
在 B 端，火山引擎的 MaaS 市场份额也远超百度智能云。
更严峻的是，传统搜索广告收入连续三个季度同比萎缩，
2025 年全年录得 162 亿元长期资产减值损失——百度正在以"左手砍旧业务、右手建新业务"
的方式完成 AI 转型，这个过程痛苦但不可逆。

\subsection{其他大厂：美团、快手、小米}
\label{subsec:china-other-bigtech}

\subsubsection{美团：无人配送的"奇点"逼近}

美团在 AI 领域的布局以"产业投资 + 场景应用"为主线。
美团龙珠（战投部门）是月之暗面天使轮和 A 轮的投资方，
同时投资了宇树科技、未来机器人、九识智能等多家具身智能企业，
瞄准仓储物流和即时配送场景。

美团从 2016 年起探索自动配送，已初步实现规模化运营：
截至 2024 年底，美团无人机已开通 53 条航线，累计配送订单超 45 万单；
自动配送车累计配送近 500 万单，自动驾驶里程占比 99\%，
全年帮助骑手减少超 240 万公里配送路程，自动驾驶总里程达 1300 万公里。
美团通过无人机（第四代新机型搭载 4D 毫米波雷达）、L4 级自动配送车"魔袋 20"
和配送机器人三条技术路线并行推进无人配送体系。
天风证券分析认为，美团无人配送正逼近"奇点"——
当单位配送成本低于人力骑手成本时，将触发指数级规模扩张。

\subsubsection{快手：可灵视频生成模型的全球影响力}

快手在多模态 AI 方面有独特积累，其可灵（Kling）视频生成模型
在第三章已有详述，是中国视频生成赛道最重要的独立玩家。
2026 年 1 月，可灵 AI 月活跃用户（MAU）突破 1200 万，
全球累计超 260 万人使用过可灵 AI，生成超 2700 万个视频和 5300 万张图片。

2026 年 1 月 31 日，可灵 3.0 系列模型正式发布并开启内测，
标志着 AI 视频创作进入新阶段。可灵 3.0 基于 All-in-One 技术理念构建，
形成多模态输入输出高度统一的一体化视频模型体系：

\begin{itemize}
  \item \textbf{视频 3.0：}单次视频生成时长最高可达 15 秒（此前为 5--10 秒），
    支持 3--15 秒灵活时长设置，引入智能分镜与自定义镜头控制；
  \item \textbf{图片 3.0：}采用视觉思维链（vCoT）技术辅助生成前的场景解构推理，
    并通过 Deep-Stack 视觉信息流机制增强细粒度感知能力；
  \item \textbf{视频 3.0 Omni：}全球首创"主体参考"功能——
    可提取 3--8 秒视频中的角色形象与音色进行还原应用，
    实现人物形象、动作与声音在复杂镜头切换中保持稳定。
\end{itemize}

多位短剧行业从业者表示，"一致性"问题一直是困扰 AI 视频生成的核心瓶颈，
可灵 3.0 在这方面取得了突破性进展，让 AI 正式进入影视与创意内容的核心生产环节。

\subsubsection{小米：从硬件公司到"AI 赋能人车家全生态"}

小米的 AI 战略与其智能硬件生态深度绑定，但 2026 年 3 月的春季发布会
标志着小米 AI 野心的质变。雷军在发布会上同时发布了三款自研大模型：

\begin{itemize}
  \item \textbf{Xiaomi MiMo-V2-Pro}：面向 Agent 时代的旗舰基座模型，
    在全球权威大模型排行榜 Artificial Analysis 上位列总榜第八、
    按品牌排名全球第五，排名超过 xAI 的 Grok；
  \item \textbf{Xiaomi MiMo-V2-Omni}：全模态大模型，
    支持文本、图像、音频、视频的统一理解与生成；
  \item \textbf{Xiaomi MiMo-V2-TTS}：语音大模型，
    为小爱同学提供更自然的对话体验。
\end{itemize}

新一代小米 SU7 首发搭载的小米 XLA 大模型，
将 HAD 辅助驾驶系统与大模型结合——用户对小爱同学说
"去那家评价最好的火锅店，帮我停在离电梯近的位置"，
AI 能实时打通大众点评评价、导航实时路况，并精准识别地库地图。
上一代 SU7 上市不到两年销售总量超 38 万辆，
2025 年全年销量 24.6 万辆，是国内 20 万以上轿车销量冠军。

雷军宣布未来三年将投入超 600 亿元于 AI 领域。
小米的独特定位在于其"物理派"路线——拥有手机（全球前三）、
汽车（SU7/YU7）、智能家居（IoT 设备超 10 亿台连接）三大硬件矩阵，
能够从多维度捕捉物理世界数据，形成"人车家全生态"的 AI 闭环。
MiMo-V2 在 OpenClaw 框架下已进化为完整的 Agent 工作流驱动核心，
能够自主推理、调用工具并作出判断。


% ============================================================
%  §11.3  INDEPENDENT STARTUP MATRIX
% ============================================================

\section{独立创业公司矩阵}
\label{sec:china-independents}

2025 年，930 家标签含"AI 应用"的中国创业公司获得新融资，
融资总额 1070.7 亿元——平均每天 2.6 家公司拿到融资，每小时 1200 万资金进场。
本节按垂直赛道逐一扫描中国 AI 独立创业公司的全貌。

\subsection{基座模型公司："六小龙"的分化与洗牌}
\label{subsec:china-foundation-models}

2023--2024 年的"百模大战"已在 2025 年完成残酷淘汰，
存活下来的头部公司被媒体称为"AI 六小龙"：
智谱、MiniMax、月之暗面（Kimi）、阶跃星辰、百川智能、零一万物。
到 2026 年初，六小龙的命运已经明显分化。

\textbf{智谱 AI}（已上市）：2021 年发布中国首个预训练大模型框架 GLM，
2024 年推出对标 GPT-4 的 GLM-4 系列。上市后市值从 579 亿港元快速攀升至约 2895 亿港元（2026 年 3 月）。
8000 家企业客户，27 万+付费开发者。OpenAI 在 2024 年 6 月的官方博客中称智谱为"头号竞争对手"。
董事长刘德兵在上市致辞中说："让机器像人一样思考，是智谱创立第一天选择的方向。"

\textbf{MiniMax}（已上市）：创始人闫俊杰，前百度实习生、前商汤副总裁，
2021 年创立时团队平均年龄仅 29 岁，董事会平均年龄 32 岁。
C 端产品（海螺 AI、Talkie）贡献超七成收入，海外收入占比超 70\%，
成为国产大模型全球化的标杆。2026 年 3 月受"龙虾热"（OpenClaw Agent 浪潮）催化，
市值两天暴涨 51\% 至 3826 亿港元，一举超越百度。
创始人闫俊杰身价超 900 亿港元，37 岁即登顶中国 AI 应用层首富。

\textbf{月之暗面/Kimi}：创始人杨植麟，清华/CMU 博士，前 Google Brain、Meta AI 研究员。
2023 年成立后以"长文本"切入，Kimi Chat 支持 20 万汉字输入。
2025 年初 DeepSeek 冲击波让 Kimi 一度陷入"至暗时刻"——
被迫收缩投放预算，经历长达 16 个月融资空窗期。
但 2025 年底发布的 K2/K2.5 模型凭借 Agent 集群能力实现反转：
K2.5 可现场调度多达 100 个分身、并行处理 1500 个步骤。
2026 年 1--3 月完成三轮融资合计超 22 亿美元，估值从 43 亿美元飙至 180 亿美元，
创下国内大模型连续融资金额最高纪录。K2.5 发布仅 20 天，
累计收入已超 2025 年全年，海外收入占比反超国内。

\textbf{阶跃星辰}：2026 年 1 月完成 50 亿元 B+ 轮融资（上海国投、国寿股权、
腾讯投资、启明创投、五源资本等），创下当月 AI 领域最大单轮融资纪录。
Step 3.5 Flash 模型在 OpenRouter 上的 Token 调用量跻身全球前列。

\textbf{百川智能}：从"AI 六小龙"中率先转向，缩减基座模型训练投入，
转向"小参数模型+行业解决方案"路线。

\textbf{零一万物}：李开复创立，同样经历了从基座模型到应用层的战略调整。

\begin{whybox}[为什么 DeepSeek 改变了全球 AI 竞争格局]
2025 年 1 月 20 日——特朗普就任第 47 任美国总统的同一天——
一家相对默默无闻的中国 AI 公司悄然发布了一个模型，数日内就撼动了全球资本市场。

DeepSeek-R1 在数学推理、编程和科学问题解决等基准测试中
匹配甚至超越了 OpenAI 的 o1 模型——然而其报告的训练成本仅约 557 万美元，
是同等美国模型投入的几十分之一。1 月 27 日，美股经历剧烈震荡：
NVIDIA 单日市值蒸发约 5930 亿美元，创下单只股票有史以来最大单日跌幅。

DeepSeek 的冲击波至少在三个层面重塑了全球 AI 竞争格局：

\textbf{第一，"算力军备竞赛"叙事被打破。}此前，行业主流观点认为训练顶级模型
必须堆砌海量 GPU，成本以亿美元计。DeepSeek 证明，通过巧妙的架构设计
（稀疏 MoE）、高效的训练策略和工程优化，
能够以极低成本达到前沿性能——这让所有基于"烧钱论"的估值逻辑受到根本质疑。

\textbf{第二，开源路线获得战略合法性。}DeepSeek-R1 以开放权重发布，
任何人都可以下载和部署。这不仅吸引了全球开发者社区的热情拥抱
（成为苹果美国 App Store 下载量第一），更迫使整个行业重新思考
"开源 vs 闭源"的战略选择——如果前沿性能可以免费获得，
闭源模型的定价权和商业壁垒将被大幅侵蚀。

\textbf{第三，出口管制的"效果悖论"再次上演。}DeepSeek 团队在芯片受限的环境下
（针对华为昇腾 910B/C 优化），反而被迫探索出了更高效的训练方法。
这与中国 AI 芯片领域的"越封锁越替代"逻辑如出一辙（详见第二章）。

进入 2026 年，DeepSeek 继续加速：V3.2 和 V3.2-Speciale 于 2025 年 11 月发布，
V4 预计将达到万亿参数规模、百万 Token 上下文窗口、原生多模态能力，
且推理成本仅约 \$0.14/M 输入 Token——约为 GPT-5 的 1/20。
Hugging Face 的分析显示，DeepSeek 和千问合计已占据全球 AI 市场约 15\% 的份额，
较一年前的 1\% 实现了数量级跃升。

然而，V4 的发布时间一再推迟，引发行业热议。
有分析认为，梁文锋团队正在攻克"长上下文一致性"这一技术难关——
在 Agent 场景中（如 OpenClaw 龙虾 Agent 需要跨越数百步推理），
模型的"健忘症"问题尤为突出。
DeepSeek 选择"宁慢不错"的发布节奏，
与字节、阿里"快速迭代、抢先发布"的风格形成鲜明对比，
也反映出 DeepSeek 作为独立研究实验室（而非大厂产品团队）的独特文化——
在这里，技术完美度优先于发布速度。

值得一提的是，DeepSeek 迄今未接受任何外部融资，
完全由幻方量化（一家量化对冲基金）自有资金支撑。
这种"不拿投资人的钱、不受资本市场节奏裹挟"的模式，
使 DeepSeek 成为全球 AI 创业生态中最独特的存在之一——
它既不像 OpenAI 那样追逐估值增长，也不像大厂 AI 团队那样服务于商业 KPI，
而是以一种近乎"学术理想主义"的方式推动 AI 前沿研究。
\end{whybox}

\subsection{具身智能：554 亿热钱涌入的超级赛道}
\label{subsec:china-embodied}

2025 年，具身智能（包括人形机器人、工业机器人、机器人零部件等）
成为中国 AI 投资最热门的单一赛道：
全年投资事件 447 起（2024 年 173 起），资本总量 554 亿元（2024 年 137 亿元），
分别增长超 250\% 和 400\%（量子位智库）。

标志性事件包括：2025 年春晚，16 台宇树机器人身着东北花袄扭秧歌，
成为全年具身智能融资狂飙的起点；银河通用机器人 B 轮 11 亿元（宁德时代领投）、
B+ 轮约 20 亿元，2026 年 3 月 A+ 轮 25 亿元刷新行业融资纪录，
估值飙至 210 亿元。

IT 桔子统计，2024--2025 年间，宇树科技、智元机器人、星动纪元、乐聚机器人、
银河通用、傅利叶、众擎机器人、云深处、星海图等公司迅速跻身独角兽行列。
具身智能已正式写入"十五五"规划，被列为国家六大未来产业之一，
同时"智能机器人"入选六大新兴支柱产业，获得政策到资本的多维度支持。

\subsubsection{头部玩家：宇树与智元的双雄对决}

进入 2026 年，中国具身智能赛道形成了宇树科技与智元机器人"双雄对决"的格局。
据 Omdia 统计，两家公司在 2025 年共同占据全球人形机器人市场 71\% 的份额。

\textbf{宇树科技}（Unitree）成立于 2016 年，总部杭州，
是全球最早实现四足机器人商业化的公司之一，占据全球四足机器人市场约 70\% 份额。
2025 年春晚出圈后热度飙升，同年 6 月完成 C 轮融资（中国移动、腾讯、
阿里巴巴、蚂蚁集团、吉利资本、锦秋资本联合领投），估值突破 100 亿元。
截至 2025 年底，累计融资超 10 亿元，投前估值超 150 亿元。
2026 年 3 月 20 日，宇树科技科创板 IPO 获上交所受理，
拟融资 42.02 亿元，目标估值约 500 亿元（约 70 亿美元），
成为 2025 年人形机器人出货量全球第一的上市候选。
创始人王兴兴年仅 35 岁，被誉为"机器人界的 DJI"。

\textbf{智元机器人}成立于 2023 年，却在不到两年内完成了从实验室到量产的飞跃。
智元的路线更偏"具身大脑"——强调 AI 驱动的通用操作能力，
在工业场景中已实现人形机器人的实际部署。
两家公司在 2026 年春晚上均有亮相：宇树展示了更精湛的武术表演，
银河通用则在沈腾、马丽身边完成盘核桃、捡玻璃碎片、叠衣理物等复杂任务。

\subsubsection{2026 年融资爆发：两个月 200 亿}

2026 年开年仅两个多月，中国具身智能赛道已披露融资超 30 起，
合计约 200 亿元——远超 2025 年一季度的 126 亿元和 2024 年一季度的 70 亿元。
这波融资潮催生了七家新晋百亿估值独角兽：
自变量机器人、星海图、灵心巧手、智平方、千寻智能、星动纪元、帕西尼感知科技。
加上此前已突破百亿估值的宇树、智元、银河通用，
中国具身智能领域已有至少 10 家百亿估值独角兽——
这一密度在全球任何单一 AI 细分赛道中都前所未见。

更值得关注的是，至少六家具身智能公司已有 2026 年内上市计划，
部分已秘密交表。上市预期是本轮融资热潮的核心动力——
投资人普遍认为，第一波上市的具身智能公司可以对照此前国产 AI 芯片公司的
上市表现获得高估值溢价。

\subsubsection{傅利叶智能与其他重要玩家}

\textbf{傅利叶智能}（Fourier Intelligence）以康复机器人起家，
逐步向通用人形机器人拓展，其 GR-2 人形机器人在 CES 2026 上引发关注。
傅利叶的差异化在于其"医疗+工业"双场景路线，
康复机器人已在全球 40 多个国家的 2000 多家医疗机构部署。

\textbf{逐际动力}获得 2 亿美元 B 轮融资，
专注于双足机器人的动态运动控制。
\textbf{灵心巧手}获近 15 亿元 B 轮融资，聚焦灵巧手操作能力。
\textbf{千寻智能}完成近 20 亿元连续两轮融资（红杉中国领投），估值突破百亿。

\subsection{AI 搜索：秘塔、天工与大厂的入口争夺}
\label{subsec:china-ai-search}

AI 搜索被认为是最具价值的 AI 应用方向之一。
国内赛道极为拥挤：秘塔 AI 搜索以"没有广告、直达结果"为卖点，
持续丰富知识库、回响（Research）等功能；天工 AI 搜索发力专业场景；
百度搜索接入深度思考模型巩固技术优势；
纳米 AI 搜索（360）集成 DeepSeek 等多家大模型；
微信搜一搜融合 DeepSeek 能力；
知乎直达升级深度思考；小红书点点深耕生活场景搜索。

从数据看，72\% 的 AI 搜索用户需求是"说问题"而非"找链接"——
这一行为变化正在根本性地重塑搜索行业。

\subsection{AI 社交与陪伴}
\label{subsec:china-ai-social}

AI 社交与陪伴是 2025 年中国 AI 应用增长最快的垂直赛道之一，
用户需求从"获取信息"延伸到"情感连接"，催生了一批差异化产品。

\textbf{MiniMax 旗下的 Talkie/海螺 AI} 是中国 AI 社交出海最成功的产品，
全球用户超 2.12 亿，以虚拟角色扮演和 AI 陪伴为核心。
Talkie 的核心竞争力在于其角色扮演模型的"情感理解"能力——
AI 角色不仅能回答问题，更能感知用户的情绪状态并主动调整对话风格，
这使其在海外市场（尤其东南亚和拉丁美洲）获得了极高的用户黏性。

\textbf{星野/猫箱（字节跳动）}基于豆包角色扮演模型构建，
主打虚拟角色互动和 AI 陪伴，凭借字节的流量分发能力快速起量。
星野的定位更偏"轻社交"——用户可以创建自己的 AI 角色并分享给好友，
形成类似"虚拟偶像+社交"的混合形态。

\textbf{自然选择（Natural Selection）}是 2026 年初异军突起的 AI 陪伴创业公司，
完成超 3000 万美元新一轮融资（阿里巴巴、蚂蚁集团、启明创投、五源资本等联合投资）。
其产品以"情绪健康"为核心差异化，将 AI 陪伴与心理健康场景结合，
切入了一个传统互联网公司难以触达的蓝海市场。

\textbf{元宝派（腾讯）}以 AI 社交新形态切入，
让用户与好友共同邀请 AI 参与聊天——这种"AI + 真实社交关系"的模式
区别于纯粹的"人与 AI 一对一"范式，利用了腾讯在社交图谱上的天然优势。

\textbf{Glow（卓明科技）}和\textbf{筑梦岛}则代表了独立创业公司在 AI 社交赛道的探索。
Glow 以"AI 分身"为核心概念，用户可以训练出具有自己说话风格的 AI 数字分身；
筑梦岛则聚焦 AI 故事创作与角色扮演社区，形成了 UGC + AI 生成的内容生态。

\subsection{AI 应用：垂直场景的崛起}
\label{subsec:china-ai-vertical-apps}

QuestMobile 2025 年度报告揭示了一个结构性转变：MAU 前十的 AI 原生 App 中，
垂直场景专业应用已占 4 席，综合助手不再是唯一的增长引擎。

\textbf{蚂蚁阿福}（蚂蚁集团）定位"AI 专业顾问"，
接入支付宝生态的金融、保险、医疗等场景知识库，
MAU 跻身国内 AI 应用 Top 5。
其差异化在于将 AI 对话能力与蚂蚁集团的金融牌照和风控体系结合，
实现"问一句话就能比价保险""AI 帮你算退税"等强场景功能。

\textbf{豆包爱学}（字节跳动）聚焦 K12 学科教育，
凭借抖音流量低成本导流快速起量，
搭载豆包大模型支持多轮对话推理，
在 QuestMobile 季度榜中稳居教育类 AI 应用前三。

\textbf{即梦 AI}（字节跳动）专注 AI 创作与设计，
以文生图、视频生成等创意工具为核心，
延续了字节"App 工厂"基因——豆包（通用）、豆包爱学（教育）、即梦（创作）
三款 App 实现"通用+垂直"精准卡位。

\textbf{秘塔 AI}以"无广告、直达结果"的 AI 搜索切入，
持续丰富知识库、回响（Research）等深度研究功能，
成为独立创业公司中 AI 搜索的标杆。
纳米 AI 搜索（360）集成 DeepSeek 等多家大模型，
小红书"点点"深耕生活场景搜索——
AI 搜索已从"找链接"演变为"说问题"，72\% 用户如此使用。

\subsection{AI 创作与设计工具}
\label{subsec:china-ai-creative}

AI 正在重构中国设计工具市场的竞争格局。
与 Figma/Adobe 主导的海外市场不同，
中国本土设计工具借助 AI 能力实现了差异化突围。

\textbf{剪映（CapCut 国内版，字节跳动）}是中国用户量最大的 AI 视频编辑工具，
集成了豆包 Seedream 生图模型和 Seedance 视频生成能力。
其"AI 一键成片""智能字幕""AI 换脸"等功能降低了短视频创作门槛，
已成为抖音内容生态的核心创作基础设施。

\textbf{美图秀秀 AI}（美图公司，港股上市）从修图工具进化为 AI 视觉平台，
AI 写真、AI 商拍、AI 视频等功能矩阵覆盖 C 端和 B 端。
2025 年 AI 相关收入占比已超 30\%，
是中国最早将 AI 能力商业化的影像公司之一。

\textbf{Pixso AI}（万兴科技）定位本土化 AI 产品设计协作平台，
主打"输入需求、AI 生成可编辑 UI 设计稿"，
支持中文语义理解和多端适配，10 万+设计师和产品经理使用。
\textbf{MasterGo AI}（莫高设计）同属国产 AI 设计赛道，
支持自然语言交互生成高保真原型与前端代码，
2025 年底用户规模快速增长，直接对标 Figma Make。
\textbf{即时设计 AI}是另一款国产在线设计协作工具，
AI 布局、AI 配色、AI 组件生成等功能
使其成为中小团队的 Figma 替代选择。

\textbf{稿定设计 AI}聚焦电商场景，
提供"AI 生成主图/详情页/海报"一站式电商视觉解决方案，
深度整合淘宝、拼多多等电商平台的设计规范。

\subsection{AI 办公与效率工具}
\label{subsec:china-ai-office}

中国 AI 办公赛道正在形成"四强争霸"格局：
WPS AI（金山办公）、飞书 AI（字节跳动）、钉钉 AI（阿里巴巴）和腾讯 AI 工作台 ima。
四家公司背后分别站着金山/小米生态、字节生态、阿里生态和腾讯生态，
AI 办公的竞争本质上是超级 App 生态之间的代理战争。

\textbf{飞书 AI/飞书妙记}（字节跳动）以 AI 原生的方式重构企业协作，
飞书妙记基于豆包大模型实现会议实时转写、智能摘要和待办提取，
并与飞书文档、飞书多维表格深度整合，
实现"开完会就有纪要、看完文档就有摘要、填完表格就有分析"的全链路 AI 体验。
飞书的差异化在于其"all-in-one"设计哲学——
将即时通讯、文档协作、日程管理、审批流程整合在单一平台内，
每个模块都嵌入了 AI 能力，避免了在多个工具间切换的效率损耗。

\textbf{腾讯 AI 工作台 ima} 月活超 1300 万，定位"AI 知识管理"，
知识库文件超 4.2 亿。ima 的独特之处在于其与微信生态的深度打通——
用户可以将微信聊天记录、公众号文章、文件等一键导入 ima 知识库，
由 AI 自动整理、分类、生成摘要，解决了"信息在微信里找不到"的高频痛点。

\textbf{WPS 灵犀}（WPS AI 3.0）于 2025 年 7 月 WAIC 上发布，
采用"左侧 Office 套件、右侧 AI 对话框"的同屏交互形态，
支持通过自然语言对话直接操作文档，被评为 WAIC"镇馆之宝"并入选全球百大 AI 应用。
截至 2025 年 6 月，WPS AI 月活用户 2951 万，WPS Office 全球月活设备 6.69 亿。
2025 年 11 月，WPS 365 升级为全球一站式 AI 协同办公平台，
推出灵犀企业版、团队空间和"轻舟"引擎，
总成本不到 Microsoft 365+Zoom+Adobe PDF 组合的三分之一。

\textbf{通义听悟}（阿里云）是国内 AI 会议转录赛道的领跑者，
基于通义千问大模型实现音视频实时转写、中英互译、智能纪要生成，
自动区分说话人并提取议程、待办和关键词。
在学习、会议、访谈等场景中积累了大量忠实用户，
与飞书妙记、讯飞听见形成三足鼎立格局。

\textbf{钉钉}在无招回归后大刀阔斧重塑产品——
2025 年推出录音卡片 DingTalk A1 和 AI 终端 DingTalk Real 两款硬件，
以"Agent OS"构建 AI 时代的工作操作系统。
钉钉的策略从叶军时代的纯平台化（PaaS only）转回"软硬一体"，
硬件+AI Agent 成为其差异化飞书的关键。

\textbf{石墨文档 AI}以"AI 数据表"和"AI 写作助手"为核心，
面向中小团队提供轻量化 AI 协作方案。
\textbf{语雀 AI}（蚂蚁集团）聚焦技术团队的知识管理，
AI 摘要、AI 搜索、AI 问答等功能使其成为开发者社区的热门选择。

\subsection{AI 编程：中国开发者工具链的崛起}
\label{subsec:china-ai-coding}

AI 编程是 2025--2026 年增长最快的开发者工具赛道之一。
全球 AI 编程工具市场由 GitHub Copilot 和 Cursor 主导，
但中国玩家正在以差异化路线快速追赶。

\textbf{通义灵码（阿里）}是国内最早推出的 AI 编程助手之一，
基于千问大模型，支持代码补全、生成、解释和单元测试编写，
已集成到 VS Code、JetBrains 等主流 IDE 中。

\textbf{Qwen Code（阿里）}于 2026 年 3 月推出，
是面向 CLI 和 IDE 的 AI 编程助手，可直接接入阿里云百炼 Coding Plan，
在千问、GLM-5、MiniMax M2.5、Kimi K2.5 等多模型间自由切换。
Qwen Code 定位对标 Claude Code，是中国首个支持多模型订阅的 AI 编程 CLI 工具。

\textbf{豆包 MarsCode（字节跳动）}是字节面向全球开发者推出的 AI 编程平台，
提供基于云端的 AI IDE 和 AI 编程助手，支持超过 100 种编程语言。
MarsCode 的差异化在于其与字节内部开发工具链的深度整合——
字节工程师使用的代码补全、代码审查、Bug 修复工具
均可通过 MarsCode 向外部开发者开放。

AI 编程赛道的竞争核心正在从"代码补全"转向"全流程 Agent"——
AI 不仅帮你写代码，还帮你理解需求、规划架构、编写测试、
部署上线甚至监控线上问题。这一趋势与 OpenClaw Agent 浪潮深度契合，
中国的 AI 编程工具正在从"辅助工具"进化为"自主编程 Agent"。

\subsection{AI 教育、AI 电商及其他垂直赛道}
\label{subsec:china-ai-verticals}

\textbf{AI 教育：}学而思旗下的 AI 学习产品持续迭代，火花 AI 聚焦数理思维，
网易有道翻译笔和学习机搭载大模型能力。科大讯飞星火大模型在教育领域的落地
最为系统化，覆盖从 K12 到成人教育的全场景——
2025 年 10 月发布"懂你的学科智能体"和"AI 学科模拟实训智能体"两款产品，
AI 学习机已迭代至第三代，首创的 AI 精准学、AI 作文批改、AI 口语陪练等功能
被同行广泛模仿。

\textbf{作业帮}依托十年教育数据积淀，形成"内容+服务+硬件"生态闭环：
直播课、会员订阅与智能学习机三大业务协同发力，
2025 年 Q3 学习机销量占比 32.6\%，稳居行业首位。
\textbf{小猿 AI}（猿辅导旗下）深耕教研领域，
标准化课程体系覆盖全学段，专职教师团队超万人，
并创新性地将心理健康守护（情绪识别技术）纳入功能体系。
AI 学习机市场整体呈爆发态势：2024 年全渠道销量 592.3 万台（+25.5\%），
销售额 190.6 亿元（+37.6\%），预计未来销售额有望突破千亿。

\textbf{AI 电商：}AI 正在重构电商的"搜索$\to$推荐$\to$决策"全链路，
从"人找货"向"AI 帮你找货、比价、下单"演进，各大平台争相布局：

\begin{itemize}
  \item \textbf{千问/淘宝问问（阿里）：}"一句话下单"功能直接调用淘宝商品信息，
    春节期间用户完成近 2 亿次下单，将 AI 对话直接打通电商交易闭环；
  \item \textbf{豆包/抖音电商（字节）：}豆包内测试 AI 购物能力——
    用户在对话框输入需求，系统直接展示抖音商城商品卡片，
    点击即可进入商品详情页完成下单和支付，无需跳转至抖音商城。
    2026 年春节期间豆包日活峰值达 1.45 亿，一旦购物功能全面开放，
    字节将拥有一个规模巨大的 AI 消费入口；
  \item \textbf{文心一言/百度优选（百度）：}以 AI 导购 + 多平台跳转为核心，
    用户通过自然语言搜索商品、比价、生成对比表，
    2026 年 3 月推出的"红手指 Operator"支持跨应用自动操作下单；
  \item \textbf{京东 AI 购：}京东探索 AI 推荐与购物助手，
    将 AI 能力整合到搜索和客服场景，深耕 3C 数码和家电等优势品类。
\end{itemize}

钛媒体分析指出，AI 改变的是购物的决策方式，但不会取代电商平台的交易根基。
当购物从"搜索商品"变成"描述需求"时，AI 有可能成为新的消费入口，
但最终的交易和履约仍然依赖于现有电商平台的供应链和物流体系。

\textbf{AI 硬件：}智能眼镜成为 2025 年最热门的 AI 硬件赛道。
全年中国市场涌现 76 款新品（同比+40\%），近百亿资金涌入。
雷鸟创新、影目科技、千问 AI 眼镜 S1（搭载阿里千问+高德/淘宝/支付宝生态）
等产品将 AI 能力从屏幕延伸到了可穿戴设备。

IDC 数据显示，2025 上半年全球智能眼镜出货 406.5 万台（+64.2\%），
其中中国市场突破百万台。SAG 统计，Rokid（乐奇）以 4\% 全球份额位列 Meta 之后第二，
华为（3\%）、小米、雷鸟创新分列三至五位——前五名中四家来自中国。
\textbf{Rokid}创始人祝铭明透露 2025 年订单超 20 万副，单 SKU 激活用户突破 10 万，
日均佩戴时长接近 8 小时。\textbf{雷鸟创新}获超 10 亿元融资（中国移动链长基金领投），
CES 2026 展示了支持 eSIM 独立联网的 X3 Pro Project。
\textbf{XREAL}与谷歌达成 Android XR 深度战略合作，
成为安卓 XR 生态的主要硬件伙伴。
\textbf{闪极科技}完成近亿元融资，专注下一代 AI 眼镜产品的定制产线。
整个赛道在 CES 2026 上由 27 家中国公司组团参展，被媒体称为"百镜大战"。

\textbf{AI 智能清洁机器人}是另一个中国占据全球领导地位的 AI 硬件品类。
2025 上半年全球出货 1535.2 万台，
\textbf{石头科技}、\textbf{科沃斯}、\textbf{追觅科技}、\textbf{云鲸}包揽全球四强，
合计份额超 54\%。中国市场前五品牌占据 89.2\% 份额，3000 元以上高端产品占比
从早期 2\% 攀升至 29\%。
追觅科技以全球首创"双足升降+底盘抬升"系统实现 8cm 越障，
AI 视觉识别 200+种障碍物；
云鲸 J5/J6 系列主打"防缠绕+智能复拖"，
石头科技 P20 Pro 接入 AI 大模型交互能力。
技术竞争已从"参数军备竞赛"转向"AI 精准解决用户痛点"——
大模型+3D 视觉+机械臂的深度结合正在重新定义智能家居。


\subsection{AI 出海：中国 App 的全球统治力}
\label{subsec:china-ai-global}

中国 AI 应用的出海成绩在 2025 年达到了历史性高度。
据 Apptopia/CNBC 统计，2025 年美国 iOS 和 Android 下载量前五的 App 中，
四款来自中国公司：TikTok、CapCut（剪映国际版）、Temu、SHEIN。
即便面临特朗普政府的持续施压和国家安全审查，中国 App 仍主导了美国下载榜。

\textbf{CapCut（剪映国际版，字节跳动）}已从视频剪辑工具进化为 AI 创作平台，
集成 AI 生图、AI 视频生成、AI 字幕翻译等功能，
是 TikTok 内容生态最重要的创作基础设施。
\textbf{TikTok}在经历近乎被禁的危机后，通过美国合资公司模式继续运营，
2025 年依然是全球下载量最大的短视频平台，
其推荐算法和 AI 内容审核系统是核心技术壁垒。

\textbf{Temu}（拼多多旗下）和\textbf{SHEIN}在 AI 驱动的供应链管理和个性化推荐方面
积累了深厚能力。SHEIN 的"小单快反"模式（AI 预测需求$\to$小批量生产$\to$快速迭代）
本质上是一套 AI 驱动的柔性供应链系统。
两者在 de minimis 贸易规则终止后仍保持美国市场增长，
展现了中国电商 AI 能力的全球竞争力。

MiniMax 旗下的 \textbf{Talkie}（全球用户超 2.12 亿）、
千问/Kimi 的海外版本，以及 DeepSeek 登顶美国 App Store 免费榜——
这些案例共同说明，中国 AI 出海已从工具类应用扩展到基座模型和 AI 原生产品层面。

字节跳动的 AI 出海版图尤为值得关注。
除了 TikTok 和 CapCut 这两款"上一代"出海产品外，
字节在 AI 原生产品层面已形成海外矩阵：
\textbf{Dola}（原 Cici，AI 助手，日活超千万）、
\textbf{Dreamina}（即梦海外版，AI 创作工具）、
\textbf{AnyGen}（工作场景 Agent，对标 Manus）。
2025 年 12 月，三款产品同步加大海外投放力度，
标志着字节 AI 出海从"单品试探"进入"矩阵作战"阶段。
36 氪评价称："C 端和 B 端的布局共同让字节具备了在海外 AI 市场参与体系化竞争的可能性。"

% ============================================================
%  §11.4  POLICY AND REGULATORY ENVIRONMENT
% ============================================================

\section{政策与监管环境}
\label{sec:china-policy}

中国对生成式 AI 的监管采取了一条"全球独特"的路径：
\textbf{备案制}——即所有向公众提供的生成式 AI 服务必须在国家网信办完成备案后方可上线。
这一制度的核心法律依据是 2023 年 8 月 15 日施行的
《生成式人工智能服务管理暂行办法》。

\subsection{大模型备案制：数据全景}
\label{subsec:china-filing-system}

截至 2025 年 12 月 31 日，累计有 748 款生成式 AI 服务完成备案，
435 款生成式 AI 应用或功能完成登记。
2025 年全年新增 446 款完成备案，平均每天超过一个服务通过审核。
截至 2026 年 2 月 28 日，这一数字进一步增长至 796 款备案、481 款登记。

从地区分布看，北京市表现突出，已完成 218 款服务备案，位居全国前列；
上海市累计完成 150 款服务备案，紧随其后。

备案制度对不同类型的 AI 服务实行分级管理：
直接开发基座模型并向公众提供服务的，需在国家网信办完成备案；
通过 API 调用已备案模型能力的应用，由地方网信办开展登记。
已上线的 AI 应用须在显著位置公示所使用的备案信息，注明模型名称和备案号。

\subsection{监管框架的多层架构}
\label{subsec:china-reg-framework}

中国 AI 监管体系呈多层架构：

\begin{itemize}
  \item \textbf{算法备案：}《互联网信息服务算法推荐管理规定》（2022 年 3 月施行），
    要求算法推荐服务提供者完成算法备案；
  \item \textbf{深度合成管理：}《互联网信息服务深度合成管理规定》（2023 年 1 月施行），
    规范 Deepfake 等技术的使用；
  \item \textbf{生成式 AI 管理：}《生成式人工智能服务管理暂行办法》（2023 年 8 月施行），
    核心制度；
  \item \textbf{数据出境安全评估：}《数据出境安全评估办法》（2022 年 9 月施行），
    对 AI 模型的跨境数据流动提出严格要求；
  \item \textbf{AI 伦理规范：}科技部发布《新一代人工智能伦理规范》，
    提出六项基本伦理要求。
\end{itemize}

\begin{warnbox}[监管不确定性：备案制的双刃剑效应]
备案制度对中国 AI 生态的影响是深刻而矛盾的。

\textbf{正面效应：}备案制建立了行业准入门槛，有效遏制了低质量模型的泛滥，
推动企业在安全性、合规性方面加大投入。对于获得备案的企业而言，
备案本身成为一种"信任凭证"，有利于获取政企客户。
此外，备案制度的"高效审核"（2025 年平均每天一个服务通过）
表明监管方并非以遏制为目的，而是试图在创新与安全之间寻求平衡。

\textbf{负面效应：}备案流程增加了创业公司的合规成本和时间成本，
对资金有限的早期团队构成不小负担。更深层的担忧在于，
备案审核中的内容安全标准具有一定的模糊性和主观性，
企业难以预判某些边界场景是否会触发监管红线，
导致"自我审查"倾向加剧，创新空间可能被压缩。

此外，备案制度与数据出境安全评估的叠加效应，
使得中国 AI 企业在全球化布局中面临独特的合规挑战——
MiniMax 海外收入占比超 70\%，其数据跨境流动的合规安排
成为上市招股书中被投资者反复追问的焦点。

\textbf{深层问题：}当监管框架的迭代速度赶不上技术发展速度时，
制度滞后本身就成为创新的最大不确定性。Agent 时代的到来
（AI 代替用户在多个应用间自主操作）对现有监管框架提出了全新挑战：
当 AI 帮你在抖音下单、在高德导航、在支付宝付款时，
责任归属和数据安全的边界在哪里？这些问题目前没有明确答案。
\end{warnbox}


% ============================================================
%  §11.5  US-CHINA COMPETITION DYNAMICS
% ============================================================

\section{中美竞争态势}
\label{sec:china-competition}

中美 AI 竞争是全球科技史上前所未有的大国技术博弈。
这场竞争的独特之处在于：双方都在 AI 的同一条赛道上全速奔跑，
但各自面临的约束条件截然不同——美国受制于成本和商业化压力，
中国受制于芯片供给和数据出境限制。
而 2025 年 1 月的 DeepSeek 冲击波，则从根本上重塑了这场竞争的叙事框架。

\subsection{DeepSeek 冲击波：一日蒸发 5930 亿美元的连锁反应}
\label{subsec:china-deepseek-shockwave}

2025 年 1 月 20 日——特朗普就任第 47 任美国总统的同一天——
DeepSeek 悄然发布 R1 模型。数日内，全球资本市场剧烈震荡：
1 月 27 日，NVIDIA 单日市值蒸发约 5930 亿美元（部分来源记录为近 6000 亿美元），
创下单只股票有史以来最大单日跌幅。
纳斯达克指数当日下跌超 3\%，费城半导体指数暴跌 9\%，
全球 AI 相关股票市值合计蒸发超过 1 万亿美元。

\begin{keybox}[DeepSeek 冲击波的全球连锁反应]
\textbf{资本市场：}NVIDIA 一日蒸发 \$5930 亿；博通、AMD、ASML 等 AI 供应链股票同步暴跌。
投资者被迫重新评估"训练顶级模型必须花费数十亿美元"的核心假设。

\textbf{行业叙事：}DeepSeek-R1 的训练成本仅约 557 万美元，
是同等性能美国模型（GPT-4 级别）投入的几十分之一。
FourWeekMBA 分析称："DeepSeek 不只是打造了一个有竞争力的 AI 模型——
它粉碎了前沿 AI 需要前沿资本的根本假设。"

\textbf{政策影响：}DeepSeek 的成功使美国出口管制面临深层悖论——
限制中国获取顶级芯片的政策，反而催生了更高效的训练方法，
证明了"约束即创新"的逆向逻辑。

\textbf{一年后的复盘：}然而，DeepSeek 发布一年后，
西方 AI 巨头不仅收复失地，更加倍下注。
NVIDIA 市值反弹至 5 万亿美元，2026 年全球 AI 领域合计资本开支达 6300 亿美元
（约占美国 GDP 的 2\%），集中在四家公司手中。
这意味着，尽管 DeepSeek 证明了效率路线的可行性，
"规模路线"并未被放弃——两种范式正在并行演进。
\end{keybox}

\subsection{"以效率换规模"vs"以规模换能力"：中美路径的根本分野}
\label{subsec:china-efficiency-vs-scale}

DeepSeek 冲击波揭示了中美 AI 竞争中一个更深层的结构性分野：
\textbf{中国走的是"以效率换规模"的路径，美国走的是"以规模换能力"的路径}。

\textbf{中国路径——效率优先：}在芯片受限的环境下，
中国公司被迫在算法层面创新，发展出稀疏 MoE 架构、更高效的训练策略和工程优化。
DeepSeek-R1 以 557 万美元训练成本达到 GPT-4 级别性能，
千问 3.5 以 3970 亿总参数仅激活 170 亿即在推理和编程任务中表现优异。
这条路径的核心逻辑是：\textit{如果你无法获得更多算力，就想办法用更少的算力做更多的事}。

\textbf{美国路径——规模优先：}OpenAI、Google、Anthropic 坚信 Scaling Laws
（规模定律）——更多数据、更大模型、更强算力必然带来更好性能。
OpenAI 的 GPT-5 训练成本预计超过 5 亿美元，
Google 2026 年 AI 资本开支预计超过 750 亿美元。
这条路径的核心逻辑是：\textit{如果你有无限算力，就用规模碾压一切}。

两条路径孰优孰劣，目前没有定论。但一个有趣的事实是：
它们正在相互学习。DeepSeek 冲击波后，美国公司开始更加重视训练效率；
而中国公司在获得更多国产芯片后，也在扩大训练规模。
最终的竞争可能不是"效率 vs 规模"的二选一，而是"效率 $\times$ 规模"的综合能力比拼。

一个耐人寻味的数据可以佐证这一判断：
DeepSeek 冲击波后一年，NVIDIA 市值不仅完全收复失地，
更进一步攀升至 5 万亿美元——这说明市场最终相信的不是"效率取代规模"，
而是"效率和规模都需要"。
与此同时，中国 AI GPU 自给率从 34\%（2024 年）预计飙升至 82\%（2027 年），
这意味着中国公司很快将同时拥有效率优势\textit{和}算力规模——
当这两条曲线在 2027--2028 年交汇时，
中美 AI 竞争将进入真正的"全面对等"阶段。

\subsection{出口管制的升级螺旋}
\label{subsec:china-export-controls}

美国对华半导体出口管制经历了四轮持续升级（详见第二章）：
2022 年 10 月限制 A100/H100 对华销售；
2023 年 10 月封堵"合规特供版"H800/A800；
2024 年 12 月将 140 家中国实体列入清单并将 HBM 纳入管制；
2026 年 3 月，特朗普政府起草全球性 AI 芯片出口许可制草案，
要求 NVIDIA、AMD 向任何国家出口先进 AI 加速器均须逐案审批。

每一轮管制的短期效果是限制中国 AI 企业获取顶级 GPU，
但中长期效果却恰好相反：华为昇腾 AI 芯片预计 2026 年占据中国市场 50\% 份额，
英伟达从 2024 年的 39\% 骤降至 8\%（Bernstein Research）。
中国 AI GPU 自给率预计从 2024 年 34\% 飙升至 2027 年 82\%（摩根士丹利），
本土 AI 芯片收入上调至 1800 亿元。

\subsection{芯片替代：从"被迫适配"到"主动选择"}
\label{subsec:china-chip-substitution}

DeepSeek V4 针对华为昇腾 910B/C 优化，成为首个完全在国产芯片上
优化的万亿参数模型。这一标志性事件意味着，中国大模型公司
已经从"被迫适配国产芯片"进入"主动为国产芯片优化"的新阶段。

更深层的变化在软件生态：华为 CANN、寒武纪 Cambricon Neuware、
摩尔线程 MUSA（CUDA 兼容层）的用户基数均在 2025 年实现数量级增长。
当客户被迫从 NVIDIA CUDA 生态迁移到国产方案后，
回迁的成本和意愿都在降低——"被迫试用"正在打破生态锁定。

阿里在 2026 年 3 月财报电话会上披露，旗下平头哥 GPU 芯片已累计规模化交付 47 万片。
百度分拆的昆仑芯于 2026 年 1 月向港交所递交 IPO 申请，目标融资 20 亿美元。
华为昇腾的客户名单已涵盖国内几乎所有头部 AI 公司——
DeepSeek V4 针对昇腾 910B/C 优化，千问和智谱均完成了昇腾适配。
这一趋势正在形成自我强化的正向循环：
越多模型适配国产芯片$\to$国产芯片软件生态越成熟$\to$更多客户愿意迁移$\to$
NVIDIA 在中国市场的份额进一步萎缩。

\subsection{开源 vs 闭源：中美路径分化}
\label{subsec:china-open-vs-closed}

中美在"开源 vs 闭源"的路线选择上出现了显著分化。

\textbf{美国：}头部公司（OpenAI、Anthropic、Google）以闭源为主，
通过 API 定价和企业订阅实现商业变现。
Meta 是例外——LLaMA 系列以开放权重发布，但动机更多是"生态锁定"而非"技术民主化"。

\textbf{中国：}开源路线获得了更广泛的采纳。DeepSeek（完全开源）、
千问（400+ 开源模型）、智谱（GLM 系列开源）、百川（部分开源）
构成了全球最活跃的开源 AI 力量之一。
Hugging Face 的分析显示，DeepSeek 和千问合计市场份额从 1\% 飙升至 15\%。

这种分化的深层原因有三：

\begin{itemize}
  \item \textbf{商业化压力差异：}美国头部公司估值以千亿美元计，
    必须通过 API 定价维持商业叙事；中国公司估值较低，
    开源反而能以最低成本获取全球开发者生态和品牌影响力；
  \item \textbf{芯片约束的倒逼：}在 GPU 受限环境下，中国公司被迫探索
    更高效的训练方法，这些方法本身就是"非硬件密集型"的创新，
    开源不会丧失核心竞争力，反而能吸引全球研究者共同优化；
  \item \textbf{政策鼓励：}中国政府在"十五五"规划和"人工智能+"行动中
    多次提及开源生态建设，国资机构也倾向于投资有开源影响力的企业。
\end{itemize}

\subsection{用户渗透率：6 亿用户重塑竞争底盘}
\label{subsec:china-user-penetration}

截至 2026 年初，中国生成式 AI 用户规模已突破 6 亿，
普及率超过四成，算力水平跃居全球前列。
QuestMobile 数据显示，AI 原生 APP 在网民中的渗透率
从 2024 年 1 月的 3\% 提升至 2024 年 12 月的 11\%，
并在 2025 年继续加速增长。

这一庞大的用户基数赋予了中国 AI 公司两个独特优势：
第一，\textbf{海量交互数据}——日均万亿级 Token 调用量产生的用户反馈，
为模型的持续优化提供了美国公司难以匹敌的数据飞轮；
第二，\textbf{极低的边际获客成本}——
当大厂通过春晚红包、免费使用等方式将 AI 助手推向亿级用户时，
独立创业公司被迫在 B 端或海外寻找差异化空间。

从全球竞争视角看，中国 AI 应用在用户规模上已形成对美国的结构性优势。
全球 MAU 前 15 的 AI 应用中，中国公司产品占据 4--5 个席位，
字节豆包在全球排名第二（仅次于 ChatGPT）。
2025 年美国 iOS 和 Android 下载量前五的 App 中，
四款来自中国公司（TikTok、CapCut、Temu、SHEIN），
加上 DeepSeek 登顶美国 App Store 免费榜——
中国 AI 产品的全球影响力已从"出海工具"扩展到"底层平台"层面。


% ============================================================
%  §11.6  MEDIA ECOSYSTEM AND INFORMATION SOURCES
% ============================================================

\section{媒体生态与信息源}
\label{sec:china-media}

理解中国 AI 创业生态，离不开一个独特的中文信息生态系统。
与硅谷依赖 TechCrunch、The Information、a16z 博客不同，
中国 AI 信息的生产和传播有自己的"基础设施"。

\subsection{核心科技媒体}
\label{subsec:china-tech-media}

\textbf{36Kr（36 氪）}是中国最具影响力的科技创投媒体。
其"暗涌 Waves"频道专注创投深度报道，"智能涌现"频道覆盖 AI 前沿。
36 氪的独家报道往往能直接影响行业认知——
本章多处引用的"豆包 DAU 破亿""千问品牌统一"等消息均首发于 36 氪。

\textbf{极客公园（GeekPark）}由张鹏创立，以深度分析和创始人访谈见长。
其年度"极客公园创新大会"是中国科技界的标志性活动。
极客公园对豆包"微笑留存曲线"的分析（180 天留存率反超 90 天留存率），
为理解 AI 产品的用户黏性提供了重要视角。

\textbf{量子位（QbitAI）}专注 AI 技术报道和行业研究，
其旗下的量子位智库发布的《2025 具身智能创投全景》等报告
是行业投资者的重要参考。2025 年 2 月的数据——
"554 亿热钱、4 大估值梯队、10 亿元现金流门槛"——出自量子位。

\textbf{晚点 LatePost}以深度长篇报道著称，
对字节跳动 AI 战略、大厂内部组织变革的报道尤为深入。
其"商业、科技与人"的定位使其在 AI 报道中保持了独特的人文视角。

\textbf{甲子光年}聚焦产业科技，以"甲子引力"大会和行业报告为核心产品，
在企业级 AI 应用报道方面有独到优势。

\subsection{数据与研究平台}
\label{subsec:china-data-platforms}

\textbf{IT 桔子}是中国最重要的创投数据库。
本章几乎所有融资数据——"2025 年 1579 起投资事件""1504.27 亿元投资总额"
"北京 405 起、深圳 263 起、上海第三"——均出自 IT 桔子。
其网站和 APP 允许按标签、轮次、地区等维度实时查询，
是中国 AI 创投研究的"基础设施级"工具。

\textbf{烯牛数据}是另一个重要的投融资数据追踪平台，
其在 2025 极新 AIGC 峰会上发布的年度报告提供了细分赛道的颗粒度分析。

\textbf{清科研究中心}隶属于清科集团，是中国最老牌的创投研究机构之一，
其年度中国股权投资排名是行业的权威参考。

\subsection{社区与即时信息}
\label{subsec:china-communities}

\textbf{即刻（Jike）}是中国 AI 从业者最活跃的社交平台。
在即刻上，创始人发产品更新、投资人分享行业观察、工程师讨论技术细节——
其在 AI 圈的地位类似于 X（Twitter）之于硅谷。
"豆包留存曲线翘尾""DeepSeek 刷屏"等话题均在即刻上首先发酵。

\textbf{投中网（China Venture）}和\textbf{投资界（PEdaily.cn）}
是创投行业的专业媒体，对融资交易、GP/LP 动态的跟踪最为及时。
本章引用的"3 年估值翻 70 倍"的具身智能分析即出自投中网。

\textbf{微信公众号生态}仍然是中国信息传播的核心渠道。
几乎所有重要的 AI 行业报告、独家消息和深度分析，
都以微信公众号文章的形式首发和传播。


% ============================================================
%  §11.7  INDUSTRY CLUSTERS
% ============================================================

\section{产业集聚：北京 vs 上海 vs 深圳 vs 杭州}
\label{sec:china-clusters}

\begin{figure}[htbp]
\centering
\begin{tikzpicture}[
  city/.style={
    circle, draw=#1, line width=1.2pt, fill=#1!10,
    align=center, font=\footnotesize,
  },
  citylbl/.style={
    font=\small\bfseries, text=#1,
  },
  tag/.style={
    font=\tiny, rounded corners=1.5pt,
    fill=#1!15, text=#1!90!black, inner sep=2pt,
  },
  stat/.style={
    font=\tiny\itshape, text=figGray,
  },
]

% --- Beijing (largest circle) ---
\node[city=figRed, minimum size=3.2cm] (bj) at (0, 0) {};
\node[citylbl=figRed] at (0, 0.8) {Beijing};
\node[stat] at (0, 0.3) {405 deals};
\node[tag=figRed] at (-0.5, -0.2) {Foundation models};
\node[tag=figRed] at (0.5, -0.7) {AI chips};

% --- Shanghai ---
\node[city=figBlue, minimum size=2.4cm] (sh) at (5.0, 0) {};
\node[citylbl=figBlue] at (5.0, 0.6) {Shanghai};
\node[stat] at (5.0, 0.15) {$\sim$250 deals};
\node[tag=figBlue] at (4.4, -0.35) {AI chips};
\node[tag=figBlue] at (5.7, -0.7) {Bio-pharma AI};

% --- Shenzhen ---
\node[city=figGreen, minimum size=2.6cm] (sz) at (10.0, 0) {};
\node[citylbl=figGreen] at (10.0, 0.7) {Shenzhen};
\node[stat] at (10.0, 0.2) {263 deals};
\node[tag=figGreen] at (9.4, -0.3) {AI hardware};
\node[tag=figGreen] at (10.7, -0.7) {Going global};

% --- Hangzhou ---
\node[city=figOrange, minimum size=2.0cm] (hz) at (14.0, 0) {};
\node[citylbl=figOrange] at (14.0, 0.5) {Hangzhou};
\node[stat] at (14.0, 0.05) {$\sim$150 deals};
\node[tag=figOrange] at (13.3, -0.4) {AI apps};
\node[tag=figOrange] at (14.6, -0.7) {E-commerce AI};

% --- Legend: circle size ---
\node[font=\tiny\itshape, text=figGray, align=center] at (7.0, -2.5)
  {Circle area $\propto$ number of AI funding deals (2025).
   Tags indicate dominant AI sub-sectors.};

% --- Representative companies below each city ---
\node[font=\tiny, text=figGray, align=center, text width=3cm] at (0, -1.9)
  {Zhipu, Moonshot AI,\\Baichuan, ByteDance};
\node[font=\tiny, text=figGray, align=center, text width=3cm] at (5.0, -1.9)
  {MiniMax, StepFun,\\Moore Threads, Biren};
\node[font=\tiny, text=figGray, align=center, text width=3cm] at (10.0, -1.9)
  {Huawei Ascend, Tencent\\Hunyuan, DJI, Horizon};
\node[font=\tiny, text=figGray, align=center, text width=3cm] at (14.0, -1.9)
  {DeepSeek, Unitree,\\Kuaishou Kling};

\end{tikzpicture}
\caption{中国 AI 创业四城地理分布对比（2025 年融资数据）}
\label{fig:china-clusters}
\end{figure}

中国 AI 产业并非均匀分布，而是高度集中于四座城市（图~\ref{fig:china-clusters}）。
IT 桔子 2025 年数据显示：
北京 405 起 AI 投融资事件（全国第一），
深圳 263 起（第二），上海紧随其后（第三，融资事件同比增速 78\%、金额翻倍），
杭州稳居第四。四座城市各有鲜明的"AI DNA"。

\begin{keybox}[北京 vs 上海 vs 深圳 vs 杭州：四城 AI DNA 对比]

\textbf{北京："模型之都"——基座模型与政策中心}

北京是中国 AI 的绝对中心。218 款大模型完成备案（全国第一），
全国超 40\% 的 AI 基座模型公司注册在北京。
智谱（清华系）、月之暗面、百川智能、零一万物总部均在北京。
百度、字节跳动（AI Lab/Seed 团队）的核心 AI 研发也在北京。
清华大学、北京大学、中科院计算所等顶尖学术资源，使北京在
AI 人才密度上独占鳌头。北京智源人工智能研究院发布的《十大 AI 技术趋势》
是行业风向标。2025 年北京 AI 获投企业数领跑全国。

\textbf{上海："芯片+金融"——AI 硬件与资本高地}

上海累计完成 150 款服务备案（全国第二），AI 投融资事件同比增速 78\%、
金额翻倍，增长势头最为强劲。上海的独特优势是"AI 芯片集聚"——
被称为"GPU 四小龙"的摩尔线程、沐曦、壁仞科技、天数智芯
全部注册在上海，集聚了全国约 40\% 的集成电路产业人才和近 50\% 的创新资源。
MiniMax（稀宇科技）总部在上海，阶跃星辰总部也在上海。
上海的另一张牌是金融——港交所和科创板的双重资本市场通道，
使上海成为 AI 公司"从融资到 IPO"的最佳跳板。
2024 年规模以上 AI 企业 394 家，产业规模 4354.92 亿元（前三季度）。

\textbf{深圳："AI+制造"——硬件落地与产业融合}

深圳以 263 起融资事件位居次席，依托其强大的电子制造业供应链，
形成了以"AI+硬件"为核心的产业特色。
华为（昇腾 AI 芯片、鸿蒙生态）、腾讯（混元/元宝核心团队）是深圳 AI 的双引擎。
大疆创新将 AI 应用于无人机视觉感知，地平线（车载 AI 芯片）在深圳有重要研发中心。
深圳的比亚迪、大族激光等制造企业与 AI 公司的深度协作，
使其成为"AI 从实验室到工厂"的最佳试验场。
深创投 2025 年布局 40 个机器人全产业链项目，国资力量在深圳的 AI 投资中尤为突出。

\textbf{杭州："产品化加速器"——从 DeepSeek 到六小龙}

杭州以"六小龙"现象（DeepSeek、宇树科技、云深处科技、游戏科学、群核科技、强脑科技）
成为 2025 年全年的科技热点。杭州 AI 投融资事件和规模稳居全国第四。
杭州的核心优势是\textbf{产品化能力}——
阿里巴巴（千问/通义）和海康威视在电商大模型与视觉大模型领域构建了极深的护城河，
服务业占 GDP 比重高达 70\% 的经济结构使杭州更擅长
依托平台与消费市场跑通"技术→产品→市场"的闭环。
杭州国资以超 10 亿元战略投资智谱，体现了地方政府"抢 AI 龙头"的决心。
高新技术企业数从 2017 年的 2844 家增至 2024 年的 1.6 万家。
\end{keybox}

四城之间正从"各自竞速"走向"协同分工"。
如果把 AI 产业理解为一条从通用能力走向落地价值的链条，
一条"源头—场景—产品"的分工正在加速成形：
\textbf{上海做源头}（基座模型 + 要素集聚 + 金融资本），
\textbf{苏州做车间}（工业场景 + 算力基建 + 规模化复制），
\textbf{杭州做产品}（消费端 + 平台生态 + 快速迭代），
\textbf{深圳做硬件}（制造业 + 供应链 + 出海通道），
\textbf{北京做大脑}（基座模型 + 政策制定 + 学术研究）。

长三角城市群的协同效应尤为突出——在 IT 桔子的 AI 投融资地域排名中，
长三角城市占据多席，苏州融资笔数同比大增 133\%，
"研发在上海、量产在苏州、市场在杭州"的区域产业闭环正在成型。

% ============================================================
%  CHAPTER SUMMARY
% ============================================================

\bigskip
\noindent\rule{\textwidth}{0.5pt}
\bigskip

回望 2025--2026 年的中国 AI 创业生态，几个核心趋势已经清晰浮现：

\textbf{第一，从"百模大战"到"寡头格局"。}基座模型赛道的淘汰赛已基本结束，
存活下来的头部公司正在资本市场兑现价值——
智谱、MiniMax 上市，月之暗面估值破千亿，阶跃星辰融资 50 亿元。
而更多的中腰部基座模型公司已被迫转向应用层或退出竞争。
"六小龙"的分化命运清晰地说明了一个残酷事实：
在基座模型赛道，只有"赢家"和"转型者"两种结局。

\textbf{第二，从软件到硬件。}具身智能成为 2025--2026 年最大的投资主题，
资本从"数字 AI"向"物理 AI"大规模迁移。
2026 年前两个月即涌入 200 亿元融资，10 家独角兽估值破百亿，
6 家公司计划年内上市。这一趋势与中国制造业供应链的天然优势高度契合——
"AI 大脑由模型公司提供，机器人身体由中国供应链制造"的分工模式正在成型。
宇树科技和智元机器人合计占据全球人形机器人市场 71\% 的份额，
这一数字或许是中国在 AI 时代最有说服力的"硬实力"指标。

\textbf{第三，大厂主导 C 端，创业公司突围 B 端和出海。}
豆包 DAU 破亿、元宝 MAU 破亿、千问 MAU 破亿——
三家大厂在 C 端的"撒币大战"已将 AI 助手的用户获取成本推至惊人高位。
独立创业公司的生存空间更多在 B 端（如智谱的政企部署）
和海外（如 MiniMax 的 Talkie、Kimi 的海外市场）。
值得关注的是，微信智能体的推出（如果成功）可能彻底改变这一格局——
14 亿月活的超级 App 接入 AI Agent，将使 C 端 AI 应用的竞争维度
从"独立 App 之争"升级为"超级 App 生态之争"。

\textbf{第四，DeepSeek 效应的深远遗产。}DeepSeek 不仅改变了全球 AI 竞争格局，
更深刻影响了中国 AI 创业的叙事逻辑——
"花大钱训大模型"不再是唯一路径，
"高效率、低成本、开源"成为新的竞争范式。
一年后的复盘显示，DeepSeek 虽然打破了"算力军备竞赛"的单一叙事，
但并未终结规模路线——美国公司 2026 年合计资本开支达 6300 亿美元，
中国公司也在加大国产芯片上的训练投入。
最终的竞争不是"效率 vs 规模"的二选一，
而是"效率 $\times$ 规模"的综合能力比拼。

\textbf{第五，AI 产品矩阵战取代单品战。}
字节跳动的案例最具代表性：豆包（通用助手）、豆包爱学（教育）、即梦（创作）、
扣子（Agent 开发）、剪映（视频编辑）、猫箱（社交陪伴）——
六款 App 共享同一个大模型底座，在不同垂直场景精准卡位。
腾讯的元宝 + ima + WorkBuddy + QClaw 矩阵、
阿里的千问 + 钉钉 + 通义听悟 + Qwen Code 矩阵，同样遵循这一逻辑。
AI 竞争正从"谁的模型最强"演变为"谁的产品生态最厚"。

\textbf{第六，监管与创新的动态平衡。}备案制度从一个侧面证明了中国
在 AI 治理方面的全球领先性——796 款服务完成备案、481 款完成登记，
这些数字本身就说明监管并非遏制创新，而是试图在安全与发展之间找到平衡点。
但随着 Agent 时代的到来，现有框架面临全新挑战：
当 AI 帮你在抖音下单、在高德导航、在支付宝付款时，
责任归属和数据安全的边界在哪里？
当微信智能体打通数百万小程序时，隐私保护将面临前所未有的压力。
中国的监管创新能否跟上技术创新的速度，是这个生态最大的悬念之一。

\textbf{第七，全球化的双向加速。}
中国 AI 出海（Cici/Dola、Talkie、DeepSeek、千问开源生态）
和海外资本入华（中东主权基金、全球 VC 追捧 IPO）正在同步加速。
中国 AI 创业生态不再是一个封闭系统，
而是全球 AI 竞争网络中不可或缺的关键节点。

中国 AI 创业生态的故事远未写完。
当"小龙虾"们在云端替人干活、当机器人在工厂里扭秧歌、
当十亿人在手机上跟 AI 聊天下单、当 6 亿用户的交互数据每天喂养着模型进化、
当宇树的人形机器人在科创板敲钟上市、当微信智能体打通 14 亿人的生活服务时——
这个全球第二大 AI 生态正在以自己的方式，
书写属于中国的 AI 创业史诗。
而这部史诗的下一章，或许将不再只是"第二极"的故事。
